\documentclass[12pt,a4paper]{article}
\usepackage[utf8]{inputenc}
\usepackage{graphicx}
\usepackage{booktabs}
\usepackage{amsmath}
\usepackage{caption}
\usepackage{float}
\usepackage{geometry}
\geometry{left=2.5cm,right=2.5cm,top=2.5cm,bottom=2.5cm}

\begin{document}

\section{SaaS Platform Performance Evaluation}
\label{sec:saas_performance}

While preceding sections authenticated functional logical boundaries, this section quantifies system efficiency under load, specifically contrasting asynchronous throughput against 12 GB edge VRAM capacities.

\subsection{Performance Topologies and Memory Load}

A concurrency stress test engaged twenty simultaneous client iterations. In purely synchronous logic, processing escalated linearly until triggering unrecoverable HTTP 504 timeouts. Conversely, the implemented asynchronous workflow (Figure \ref{fig:async_vs_sync}) guaranteed HTTP 202 reception beneath 68 milliseconds universally, successfully bounding system throughput exclusively to GPU limits (124 images/min) without I/O thread collapse.

\begin{figure}[htbp]
    \centering
    \begin{minipage}{0.48\textwidth}
        \centering
        \includegraphics[width=\textwidth]{figures/plot_async_vs_sync.png}
        \caption{Asynchronous throughput resilience.}
        \label{fig:async_vs_sync}
    \end{minipage}\hfill
    \begin{minipage}{0.48\textwidth}
        \centering
        \includegraphics[width=\textwidth]{figures/plot_memory_stability.png}
        \caption{Registry lazy-load VRAM stability.}
        \label{fig:memory_stability}
    \end{minipage}
\end{figure}

Simultaneously, booting multiple deep models categorically provokes hardware exhaustion. Employing the Lazy-Load registry strictly capped allocation (Figure \ref{fig:memory_stability}). Despite Mask R-CNN consuming 5.5GB and YOLOv8s 1.5GB, long-term switching across mixed loads exhibited 0GB memory leakage, stabilising perfectly at 7.4GB against the 12GB ceiling constraint.

\subsection{Execution Latency Breakdown}

Granular hardware time-probes isolated the sub-routines comprising the Infrared hybrid pipeline. As shown, the purely algorithmic pixel-to-temperature component—the central contribution—absorbs only 10\% of the load, proving hyper-efficient computational design.

\begin{table}[htbp]
    \centering
    \caption{Empirical Hardware Latency Decomposition}
    \vspace{0.2cm}
    \begin{tabular}{@{}lcc@{}}
        \toprule
        \textbf{Operation Stage} & \textbf{Time (ms)} & \textbf{Proportion} \\
        \midrule
        OCR Temp Extraction (CPU) & \char`~ 520 ms & 68\% \\
        Algorithmic Pixel Mapping (GPU/CPU) & \char`~ 80 ms & 10\% \\
        Topological Hotspot Tracking & \char`~ 50 ms & 7\% \\
        Mask I/O Generation & \char`~ 115 ms & 15\% \\
        \midrule
        \textbf{Comprehensive Pipeline Delay} & \textbf{\char`~ 765 ms} & \textbf{100\%} \\
        \bottomrule
    \end{tabular}
    \label{tab:latency_breakdown}
\end{table}

Ultimately, the SaaS integration satisfies every performance objective: processing over 120 images per minute concurrently without cross-tenant breaches, safely bounded within VRAM targets.

\end{document}
